\section*{Kurze Einsteigerübersicht: Excel benutzen}

\begin{boxK}
\textbf{Worum geht es?} Excel ist ein Tabellenprogramm. Du kannst Daten eintragen, sortieren, mit Formeln rechnen und Ergebnisse als Diagramm darstellen. Wichtig ist: Excel rechnet nur das, was du als Befehl eingibst.
\end{boxK}

\subsection*{1. Die wichtigsten Begriffe}
\renewcommand{\arraystretch}{1.35}
\begin{tabularx}{\textwidth}{|>{\raggedright\arraybackslash}p{3.0cm}|>{\raggedright\arraybackslash}X|>{\raggedright\arraybackslash}p{3.0cm}|}
\hline
\rowcolor{gray!20}\textbf{Begriff} & \textbf{Bedeutung} & \textbf{Beispiel} \\\hline
Arbeitsmappe & die ganze Excel-Datei & Schulfest-Kiosk.xlsx \\\hline
Tabellenblatt & eine einzelne Seite in der Datei & Blatt 1, Messreihe \\\hline
Spalte & senkrechter Bereich & A, B, C, D \\\hline
Zeile & waagerechter Bereich & 1, 2, 3, 4 \\\hline
Zelle & Schnittpunkt aus Spalte und Zeile & C3 \\\hline
Zellbezug & Verweis auf eine Zelle & C3 bedeutet: Wert aus C3 \\\hline
\end{tabularx}

\subsection*{2. Daten richtig eintragen}
\begin{boxK}
\textbf{Regel:} Lies zuerst die Zeile, danach die Spalte. Trage den Wert dort ein, wo sich Zeile und Spalte treffen.
\end{boxK}

\textbf{Beispiel:} Der Messwert für \textbf{Zeile 1, Spalte 3} lautet \textbf{12}. Dann gehört die Zahl 12 in das Feld bei Zeile 1 und Spalte 3.

\vspace{0.2cm}
\begin{center}
\begin{tabular}{|c|c|c|c|}
\hline
 & \textbf{Spalte 1} & \textbf{Spalte 2} & \textbf{Spalte 3} \\\hline
\textbf{Zeile 1} & & & 12 \\\hline
\textbf{Zeile 2} & & & \\\hline
\textbf{Zeile 3} & & & \\\hline
\end{tabular}
\end{center}

\newpage
\subsection*{3. Formeln verstehen}
\begin{boxK}
\textbf{Merke:} Jede Formel beginnt mit einem Gleichheitszeichen. Ohne \texttt{=} behandelt Excel die Eingabe wie normalen Text.
\end{boxK}

\renewcommand{\arraystretch}{1.35}
\begin{tabularx}{\textwidth}{|>{\raggedright\arraybackslash}p{3.2cm}|>{\raggedright\arraybackslash}p{4.2cm}|>{\raggedright\arraybackslash}X|}
\hline
\rowcolor{gray!20}\textbf{Befehl / Zeichen} & \textbf{Beispiel} & \textbf{Bedeutung} \\\hline
\texttt{=} & \texttt{=C3*D3} & Excel soll rechnen. \\\hline
\texttt{*} & \texttt{=C3*D3} & Wert aus C3 mal Wert aus D3. \\\hline
\texttt{-} & \texttt{=E3-F3} & Wert aus F3 wird von E3 abgezogen. \\\hline
\texttt{SUMME} & \texttt{=SUMME(G3:G8)} & Alle Werte von G3 bis G8 werden addiert. \\\hline
\texttt{MITTELWERT} & \texttt{=MITTELWERT(G3:G8)} & Der Durchschnitt der Werte von G3 bis G8 wird berechnet. \\\hline
\texttt{\$} & \texttt{\$B\$13} & Diese Zelle bleibt beim Herunterziehen fest. \\\hline
\end{tabularx}

\section*{So gehst du bei einer Excel-Aufgabe vor}

\begin{boxK}
\textbf{Drei Fragen vor jeder Formel:}
\begin{enumerate}[label=\arabic*)]
  \item Welche Werte brauche ich?
  \item In welchen Zellen stehen diese Werte?
  \item Welche Rechenoperation passt dazu?
\end{enumerate}
\end{boxK}

\subsection*{4. Beispiel: Umsatz berechnen}
\begin{tabularx}{\textwidth}{|>{\raggedright\arraybackslash}p{5.2cm}|>{\raggedright\arraybackslash}X|}
\hline
\rowcolor{gray!20}\textbf{Frage} & \textbf{Antwort} \\\hline
Welche Werte brauche ich? & Verkaufspreis und verkaufte Anzahl \\\hline
Wo stehen die Werte? & Verkaufspreis in C3, Anzahl in D3 \\\hline
Welche Rechnung passt? & Verkaufspreis mal Anzahl \\\hline
Welche Excel-Formel passt? & \texttt{=C3*D3} \\\hline
\end{tabularx}

\subsection*{5. Ergebnisse kontrollieren}
\begin{boxK}
\textbf{Kontrolliere nicht nur, ob Excel eine Zahl ausgibt. Kontrolliere, ob die Zahl sinnvoll ist.}
\end{boxK}

\begin{itemize}
  \item Ist die Formel mit \texttt{=} gestartet?
  \item Werden die richtigen Zellen verwendet?
  \item Passt die Rechenart? Zum Beispiel: Umsatz = Preis mal Anzahl.
  \item Ist das Ergebnis ungefähr plausibel?
  \item Wurde die Formel richtig nach unten übertragen?
\end{itemize}

\subsection*{6. Mini-Check zum Ausfüllen}
\begin{enumerate}[label=\alph*)]
  \item Was bedeutet die Zelle \textbf{D5}?\\[0.6cm]
  \rule{\textwidth}{0.4pt}

  \item Was berechnet die Formel \texttt{=C3*D3}?\\[0.6cm]
  \rule{\textwidth}{0.4pt}

  \item Was bedeutet der Bereich \texttt{G3:G8}?\\[0.6cm]
  \rule{\textwidth}{0.4pt}

  \item Warum ist \texttt{=SUMME(G3:G8)} oft besser als alle Werte einzeln einzutippen?\\[0.6cm]
  \rule{\textwidth}{0.4pt}
\end{enumerate}

\begin{boxK}
\textbf{Merksatz:} Eine gute Excel-Lösung besteht nicht nur aus richtigen Zahlen. Du solltest erklären können, welche Zellen deine Formel benutzt und warum diese Rechnung zum Problem passt.
\end{boxK}
